% !TEX root = 第七章_二次型(答案版).tex
\setcounter{chapter}{6}
\pagestyle{fancy}

\chapter{二次型}

\section{双线性型}

\begin{definition}
设 $V',V,W$ 是域 $F$ 上的线性空间.若映射
\[
g:V'\times V\to W
\]
对两个变量分别都是线性的,则称 $g$ 为一个双线性型.常记
\[
g(\alpha,\beta)=\langle\alpha,\beta\rangle.
\]
\end{definition}

若 $V'$ 与 $V$ 分别取基
\[
\varepsilon_1,\ldots,\varepsilon_m,
\qquad
\eta_1,\ldots,\eta_n,
\]
且
\[
\alpha=(\varepsilon_1,\ldots,\varepsilon_m)x,
\qquad
\beta=(\eta_1,\ldots,\eta_n)y,
\]
则
\[
g(\alpha,\beta)=x^{\mathsf T}Gy,
\qquad
G=(g_{ij}),
\qquad
g_{ij}=g(\varepsilon_i,\eta_j).
\]

\begin{theorem}[\markstar]
固定 $V',V$ 的基以后,$V'\times V$ 上的双线性型与 $F$ 上的
$m\times n$ 矩阵一一对应.若分别换基
\[
x=P\widehat x,\qquad y=Q\widehat y,
\]
则双线性型的矩阵由 $G$ 变为
\[
\widehat G=P^{\mathsf T}GQ.
\]
\end{theorem}

\begin{definition}
若存在非零 $\alpha\in V'$,使
\[
g(\alpha,\beta)=0
\qquad
(\forall\,\beta\in V),
\]
则称 $g$ 左退化;右退化类似定义.若既非左退化也非右退化,则称
$g$ 非退化.
\end{definition}

由矩阵表示立即得到:
\[
g\text{ 非右退化}\iff r(G)=n,
\]
\[
g\text{ 非左退化}\iff r(G)=m.
\]
有限维情况下,
\[
g\text{ 非退化}
\iff
m=n=r(G)
\iff
G\text{ 可逆}.
\]

\begin{theorem}[\markstar]
设 $V',V$ 都有限维,则 $V'\times V$ 上双线性型 $g$ 非退化的充分必要条件为
\[
\dim V'=\dim V
\]
且其表示矩阵可逆.
\end{theorem}

\begin{definition}
对线性映射
\[
\varphi:V_1\to V_2,
\]
定义对偶映射
\[
\varphi^*:V_2^*\to V_1^*,
\qquad
\varphi^*(f)=f\circ\varphi.
\]
\end{definition}

\begin{theorem}[\markstar]
若 $\varphi$ 在 $V_1,V_2$ 的基下矩阵为 $A$,而
$\varphi^*$ 在相应对偶基下矩阵为 $B$,则
\[
\boxed{B=A^{\mathsf T}}.
\]
\end{theorem}

下面取 $V'=V$.

\begin{definition}
若
\[
g(\alpha,\beta)=g(\beta,\alpha)
\]
对任意 $\alpha,\beta\in V$ 成立,则称 $g$ 为对称双线性型.
在给定基下,其矩阵 $G$ 为对称矩阵.相应函数
\[
Q(\alpha)=g(\alpha,\alpha)
\]
称为二次型.
\end{definition}

在特征不为 $2$ 的数域上,二次型与对称双线性型可由极化公式互相确定:
\[
g(\alpha,\beta)
=
\dfrac12
\left(
Q(\alpha+\beta)-Q(\alpha)-Q(\beta)
\right).
\]

若实对称双线性型满足
\[
g(\alpha,\alpha)>0
\qquad
(\forall\,\alpha\neq0),
\]
则称其为正定双线性型;对应二次型也称为正定二次型.


%例题7.1.1
\begin{example}
设 $\mathbb{R}[x]$ 为实多项式空间.对 $f,g\in\mathbb{R}[x]$ 定义
\[
\varphi(f,g)=f(x)g'(x).
\]
证明 $\varphi$ 是 $\mathbb{R}[x]$ 上的双线性映射.
\end{example}
\exampleblank

\begin{solution}
对固定 $g$,
\[
\varphi(af_1+bf_2,g)
=
(af_1+bf_2)g'
=
a\varphi(f_1,g)+b\varphi(f_2,g).
\]
又因求导是线性映射,对固定 $f$,
\[
\varphi(f,ag_1+bg_2)
=
f(ag_1'+bg_2')
=
a\varphi(f,g_1)+b\varphi(f,g_2).
\]
因此 $\varphi$ 对两个变量分别线性.
\end{solution}


%例题7.1.2
\begin{example}
设
\[
A=(a_{ij})_{2\times2}\in M_2(K),
\]
在 $M_2(K)$ 上定义
\[
f(X,Y)=\tr(X^{\mathsf T}AY).
\]
\begin{enumerate}
    \item 证明 $f$ 是双线性型;
    \item 求它在基
    \[
    E_{11},E_{12},E_{21},E_{22}
    \]
    下的矩阵;
    \item 判断 $f$ 何时非退化.
\end{enumerate}
\end{example}
\exampleblank

\begin{solution}
迹,矩阵乘法以及转置对相应变量均为线性的,所以第一问成立.

按基顺序
\[
E_{11},E_{12},E_{21},E_{22}
\]
计算
\[
G_{rs}=f(E_r,E_s),
\]
得到
\[
\boxed{
G=
\begin{pmatrix}
a_{11}&0&a_{12}&0\\
0&a_{11}&0&a_{12}\\
a_{21}&0&a_{22}&0\\
0&a_{21}&0&a_{22}
\end{pmatrix}
=
A\otimes I_2.
}
\]
因此
\[
\det G=(\det A)^2.
\]
故
\[
\boxed{f\text{ 非退化}\iff \det A\neq0.}
\]
\end{solution}


%例题7.1.3
\begin{example}
证明 $M_n(K)$ 上的双线性型
\[
f(X,Y)=\tr(XY)
\]
是非退化的.
\end{example}
\exampleblank

\begin{solution}
设 $X\neq0$,取某个 $x_{ij}\neq0$.令
\[
Y=E_{ji}.
\]
则
\[
\tr(XY)=x_{ij}\neq0.
\]
因此不存在非零 $X$ 与全部 $Y$ 正交;右侧同理,故该双线性型非退化.
\end{solution}


%例题7.1.4
\begin{example}
设 $V$ 是数域 $K$ 上的 $n$ 维线性空间,
$f(\alpha,\beta)$ 是 $V$ 上的非退化双线性型.证明:对任意
$g\in V^*$,存在唯一 $\alpha\in V$,使
\[
g(\beta)=f(\alpha,\beta)
\qquad
(\forall\,\beta\in V).
\]
\end{example}
\exampleblank

\begin{solution}
定义线性映射
\[
T:V\to V^*,
\qquad
T(\alpha)=f(\alpha,\cdot).
\]
非退化性说明
\[
\ker T=\{0\},
\]
故 $T$ 单射.又
\[
\dim V=\dim V^*=n,
\]
所以 $T$ 是同构.于是任意 $g\in V^*$ 都有唯一原像 $\alpha$.
\end{solution}


%例题7.1.5
\begin{example}
设 $f(\alpha,\beta)$ 是线性空间 $V$ 上的对称或反对称双线性型,
$W$ 是 $V$ 的真子空间.证明:对任意 $\xi\notin W$,存在非零
\[
\eta\in W+\Span\{\xi\}
\]
使
\[
f(\eta,\alpha)=0
\qquad
(\forall\,\alpha\in W).
\]
\end{example}
\exampleblank

\begin{solution}
令
\[
U=W+\Span\{\xi\}.
\]
因为 $\xi\notin W$,
\[
\dim U=\dim W+1.
\]
定义
\[
T:U\to W^*,
\qquad
T(\eta)=f(\eta,\cdot)|_W.
\]
有
\[
\dim U>\dim W^*,
\]
故 $\ker T\neq\{0\}$.取非零 $\eta\in\ker T$ 即得结论.
\end{solution}


%例题7.1.6
\begin{example}
设 $f(\alpha,\beta)$ 是数域 $K$ 上线性空间 $V$ 上的双线性型,定义
\[
V_1=\{x\in V:f(x,y)=0,\ \forall y\in V\},
\]
\[
V_2=\{y\in V:f(x,y)=0,\ \forall x\in V\}.
\]
证明 $V_1,V_2$ 都是 $V$ 的子空间.
\end{example}
\exampleblank

\begin{solution}
以 $V_1$ 为例.若 $x_1,x_2\in V_1$,$a,b\in K$,则对任意 $y$,
\[
f(ax_1+bx_2,y)
=
af(x_1,y)+bf(x_2,y)=0.
\]
故 $ax_1+bx_2\in V_1$.$V_2$ 同理.
\end{solution}


%例题7.1.7
\begin{example}
设 $V',V$ 是域 $F$ 上的 $n$ 维线性空间,
\[
g:V'\times V\to F
\]
为双线性型.证明下列条件等价:
\begin{enumerate}
    \item $g$ 非左退化;
    \item $g$ 非右退化;
    \item $g$ 非退化;
    \item $g$ 在任意基下的矩阵均可逆.
\end{enumerate}
\end{example}
\exampleblank

\begin{solution}
取任意两组基,设 $g$ 的矩阵为 $G\in M_n(F)$.
有
\[
g\text{ 非左退化}\iff r(G)=n,
\]
\[
g\text{ 非右退化}\iff r(G)=n.
\]
因此前三条等价,并且均等价于 $G$ 可逆.

换基后矩阵为
\[
P^{\mathsf T}GQ,
\]
其中 $P,Q$ 可逆,所以可逆性与基的选择无关,第四条也等价.
\end{solution}


%例题7.1.8
\begin{example}
设
\[
g:V'\times V\to F
\]
为双线性型,定义左核,右核
\[
k_l(g)=\{\alpha\in V':g(\alpha,\beta)=0,\ \forall\beta\in V\},
\]
\[
k_r(g)=\{\beta\in V:g(\alpha,\beta)=0,\ \forall\alpha\in V'\}.
\]
证明:
\begin{enumerate}
    \item $k_l(g),k_r(g)$ 都是子空间;
    \item $g$ 自然诱导出
    \[
    \overline g:(V'/k_l(g))\times V\to F,
    \]
    且 $\overline g$ 非左退化;
    \item 进一步诱导出
    \[
    \widetilde g:(V'/k_l(g))\times(V/k_r(g))\to F,
    \]
    且 $\widetilde g$ 非退化.
\end{enumerate}
\end{example}
\exampleblank

\begin{solution}
第一问由双线性直接验证.

定义
\[
\overline g(\alpha+k_l(g),\beta)=g(\alpha,\beta).
\]
若更换 $\alpha$ 的陪集代表元,只会增加左核中的向量,故定义良好.
若某陪集对所有 $\beta$ 都取零,则其代表元属于左核,所以该陪集为零,
故 $\overline g$ 非左退化.

再定义
\[
\widetilde g(\alpha+k_l(g),\beta+k_r(g))=g(\alpha,\beta).
\]
同理定义良好.若任一侧元素与另一侧全部元素配对均为零,则其代表元分别落入
对应核中,因此两侧均非退化.
\end{solution}


%例题7.1.9
\begin{example}[\markstar]
设 $f(\alpha,\beta)$ 是实数域上 $n$ 维线性空间 $V$ 上的对称双线性型,
其在某基下的矩阵为 $A$.已知二次型
\[
q(x)=x^{\mathsf T}Ax
\]
的负惯性指数为 $0$.令
\[
W=\{\alpha\in V:f(\alpha,\alpha)=0\}.
\]
证明 $W$ 是子空间,并求 $\dim W$.
\end{example}
\exampleblank

\begin{solution}
负惯性指数为 $0$,故 $A$ 半正定.半正定矩阵满足
\[
x^{\mathsf T}Ax=0
\iff
Ax=0.
\]
因此
\[
W=\ker A.
\]
故 $W$ 是子空间,并且
\[
\boxed{\dim W=n-r(A)}.
\]
\end{solution}


%例题7.1.10
\begin{example}
设
\[
f(x,y)=x^{\mathsf T}Ay,
\qquad
x,y\in\mathbb{R}^{2n+1},
\]
其中
\[
A=
\begin{pmatrix}
1&0&0\\
0&O&I_n\\
0&I_n&O
\end{pmatrix}.
\]
令
\[
C=
\{X\in M_{2n+1}(\mathbb{R}):
f(Xu,v)=-f(u,Xv),\ \forall u,v\}.
\]
证明 $C$ 是维数
\[
2n^2+n
\]
的线性子空间,并给出一组基.
\end{example}
\exampleblank

\begin{solution}
条件等价于
\[
X^{\mathsf T}A+AX=0.
\]
按 $1,n,n$ 分块写
\[
X=
\begin{pmatrix}
a&p^{\mathsf T}&q^{\mathsf T}\\
r&B&C_1\\
s&D&E
\end{pmatrix}.
\]
代入得到
\[
a=0,\qquad s=-p,\qquad r=-q,
\]
\[
D^{\mathsf T}=-D,\qquad C_1^{\mathsf T}=-C_1,\qquad E=-B^{\mathsf T}.
\]
因此
\[
\boxed{
X=
\begin{pmatrix}
0&p^{\mathsf T}&q^{\mathsf T}\\
-q&B&C_1\\
-p&D&-B^{\mathsf T}
\end{pmatrix},
}
\]
其中 $p,q\in\mathbb{R}^n$,$B$ 任意,而 $C_1,D$ 反对称.

自由参数个数为
\[
2n+n^2+2\cdot\dfrac{n(n-1)}2
=
\boxed{2n^2+n}.
\]

一组基可由下列参数逐个取为标准基得到:
\[
p=e_i,\quad q=e_i,\quad B=E_{ij},
\]
以及 $C_1,D$ 分别取
\[
E_{ij}-E_{ji}\qquad(i<j),
\]
其余参数取零.
\end{solution}


%例题7.1.11
\begin{example}
设 $f(\alpha,\beta)$ 是数域 $K$ 上线性空间 $V$ 上的对称双线性型,
并且存在 $f_1,f_2\in V^*$ 使
\[
f(\alpha,\beta)=f_1(\alpha)f_2(\beta).
\]
证明:存在非零常数 $k\in K$ 与线性函数 $g\in V^*$,使
\[
f(\alpha,\beta)=k\,g(\alpha)g(\beta).
\]
\end{example}
\exampleblank

\begin{solution}
若 $f\equiv0$,取任意 $k\neq0$ 与 $g=0$ 即可.

以下设 $f\neq0$.对称性给出
\[
f_1(\alpha)f_2(\beta)
=
f_1(\beta)f_2(\alpha).
\]
取 $\alpha_0$ 使 $f_1(\alpha_0)\neq0$,则对任意 $\beta$,
\[
f_2(\beta)
=
\dfrac{f_2(\alpha_0)}{f_1(\alpha_0)}f_1(\beta).
\]
故 $f_2=c f_1$,其中 $c\neq0$.令
\[
g=f_1,\qquad k=c,
\]
即得结论.
\end{solution}


%例题7.1.12
\begin{example}
\begin{enumerate}
    \item 设 $V$ 是实线性空间,$f$ 是 $V$ 上正定对称双线性型,
    $U$ 是有限维子空间.证明
    \[
    V=U\oplus U^\perp.
    \]
    \item 设 $V$ 是数域 $K$ 上的 $n$ 维线性空间,$g$ 是非退化对称双线性型,
    $W$ 是子空间.证明
    \[
    \dim V=\dim W+\dim W^\perp,
    \qquad
    (W^\perp)^\perp=W.
    \]
\end{enumerate}
\end{example}
\exampleblank

\begin{solution}
第一问取 $U$ 的一组基并对 $f$ 作 Gram--Schmidt 正交化,
得到正交基 $e_1,\ldots,e_r$.对任意 $v\in V$,令
\[
u=
\sum\limits_{i=1}^{r}
\dfrac{f(v,e_i)}{f(e_i,e_i)}e_i.
\]
则 $u\in U$ 且 $v-u\in U^\perp$.又正定性保证
\[
U\cap U^\perp=\{0\},
\]
所以为直和.

第二问定义
\[
T:V\to W^*,
\qquad
T(v)=g(v,\cdot)|_W.
\]
非退化性保证 $T$ 满射,而
\[
\ker T=W^\perp.
\]
故
\[
n=\dim W+\dim W^\perp.
\]
显然
\[
W\subseteq(W^\perp)^\perp,
\]
两者维数相同,故相等.
\end{solution}


\newpage

\section{二次型的标准形}

\begin{theorem}[\markstar]
Witt 消去定理给出合同分解中的消去原则;在实数域上,惯性定理说明
任意实二次型都可经非退化实线性替换化为
\[
y_1^2+\cdots+y_p^2-y_{p+1}^2-\cdots-y_r^2,
\]
其中正惯性指数 $p$,负惯性指数 $r-p$ 唯一确定.
\end{theorem}


\subsection{基础计算}

%例题7.2.1
\begin{example}
用可逆线性代换化下列二次型为规范形,并说明其惯性指数:
\begin{enumerate}
    \item
    \[
    4x_1^2+x_2^2+x_3^2-4x_1x_2+4x_1x_3-3x_2x_3;
    \]
    \item
    \[
    x_1^2+5x_2^2-4x_3^2+2x_1x_2-4x_2x_3;
    \]
    \item
    \[
    x_1^2-3x_2^2-2x_1x_2+2x_1x_3-6x_2x_3;
    \]
    \item
    \[
    x_1x_{2n}+x_2x_{2n-1}+\cdots+x_nx_{n+1};
    \]
    \item
    \[
    x_1x_2+x_2x_3+\cdots+x_{n-1}x_n;
    \]
    \item
    \[
    \displaystyle\sum\limits_{i=1}^{n}x_i^2+
    \sum\limits_{1\le i,j\le n}x_ix_j;
    \]
    \item
    \[
    \displaystyle\sum\limits_{i=1}^{n}(x_i-\overline x)^2,
    \qquad
    \overline x=\dfrac{x_1+\cdots+x_n}{n}.
    \]
\end{enumerate}
\end{example}
\exampleblank

\begin{solution}
\begin{enumerate}
    \item 令
    \[
    y_1=2x_1-x_2+x_3,\qquad
    y_2=\dfrac{x_2-x_3}{2},\qquad
    y_3=\dfrac{x_2+x_3}{2}.
    \]
    则
    \[
    Q=y_1^2+y_2^2-y_3^2.
    \]
    故正,负惯性指数为 $2,1$.

    \item 先写成
    \[
    Q=(x_1+x_2)^2+(2x_2-x_3)^2-5x_3^2.
    \]
    再令 $y_3=\sqrt5\,x_3$,即得
    \[
    Q=y_1^2+y_2^2-y_3^2.
    \]
    正,负惯性指数为 $2,1$.

    \item
    \[
    Q=(x_1-x_2+x_3)^2-(2x_2+x_3)^2.
    \]
    取第三个独立新变量补成可逆变换,得到规范形
    \[
    Q=y_1^2-y_2^2.
    \]
    正,负,零指数分别为 $1,1,1$.

    \item 对每一对
    \[
    (x_i,x_{2n+1-i})
    \]
    作代换
    \[
    u_i=\dfrac{x_i+x_{2n+1-i}}{2},
    \qquad
    v_i=\dfrac{x_i-x_{2n+1-i}}{2},
    \]
    则
    \[
    x_ix_{2n+1-i}=u_i^2-v_i^2.
    \]
    因而规范形有 $n$ 个正平方与 $n$ 个负平方.

    \item 该二次型的对称矩阵只有相邻两条副对角线上元素为 $\dfrac12$.
    其特征值为
    \[
    \lambda_k=\cos\dfrac{k\pi}{n+1},
    \qquad k=1,\ldots,n.
    \]
    对应正交特征向量可取
    \[
    \left(
    \sin\dfrac{k\pi}{n+1},
    \sin\dfrac{2k\pi}{n+1},
    \ldots,
    \sin\dfrac{nk\pi}{n+1}
    \right)^{\mathsf T}.
    \]
    因此正,负惯性指数均为
    \[
    \left\lfloor\dfrac n2\right\rfloor,
    \]
    当 $n$ 为奇数时另有一个零平方项.

    \item
    \[
    Q=\|x\|^2+(x_1+\cdots+x_n)^2.
    \]
    在向量 $(1,\ldots,1)^{\mathsf T}$ 方向上的特征值为 $n+1$,
    在其正交补上的特征值均为 $1$,故为正定二次型,规范形为
    \[
    y_1^2+\cdots+y_n^2.
    \]

    \item
    \[
    Q=
    \sum\limits_i x_i^2
    -
    \dfrac1n
    \left(\sum\limits_i x_i\right)^2.
    \]
    对称矩阵为
    \[
    I-\dfrac1nJ.
    \]
    在 $(1,\ldots,1)^{\mathsf T}$ 上特征值为 $0$,其正交补上特征值为 $1$.
    因此规范形为
    \[
    y_1^2+\cdots+y_{n-1}^2,
    \]
    秩为 $n-1$.
\end{enumerate}
\end{solution}


%例题7.2.2
\begin{example}
设 $A$ 为 $n$ 阶可逆实矩阵,
\[
B=
\begin{pmatrix}
O&A\\
A^{\mathsf T}&O
\end{pmatrix}.
\]
求 $B$ 的正,负惯性指数.
\end{example}
\exampleblank

\begin{solution}
作合同变换可先把 $A$ 化成 $I_n$,于是 $B$ 合同于
\[
\begin{pmatrix}
O&I_n\\
I_n&O
\end{pmatrix}.
\]
再用
\[
u=\dfrac{x+y}{\sqrt2},
\qquad
v=\dfrac{x-y}{\sqrt2},
\]
可化为
\[
\diag(I_n,-I_n).
\]
因此
\[
\boxed{p(B)=q(B)=n}.
\]
\end{solution}


%例题7.2.3
\begin{example}
设
\[
f(x)=
\sum\limits_{i=1}^{s}
(a_{i1}x_1+\cdots+a_{in}x_n)^2.
\]
证明 $f$ 的秩等于矩阵
\[
A=
\begin{pmatrix}
a_{11}&\cdots&a_{1n}\\
\vdots&&\vdots\\
a_{s1}&\cdots&a_{sn}
\end{pmatrix}
\]
的秩.
\end{example}
\exampleblank

\begin{solution}
有
\[
f(x)=\|Ax\|^2=x^{\mathsf T}A^{\mathsf T}Ax.
\]
且
\[
\ker(A^{\mathsf T}A)=\ker A,
\]
故
\[
r(A^{\mathsf T}A)=r(A).
\]
二次型的秩即其对称矩阵 $A^{\mathsf T}A$ 的秩,所以结论成立.
\end{solution}


%例题7.2.4
\begin{example}
设
\[
f(x)=x^{\mathsf T}Ax
\]
是实二次型,并存在 $x_1,x_2\in\mathbb{R}^n$ 使
\[
x_1^{\mathsf T}Ax_1>0,
\qquad
x_2^{\mathsf T}Ax_2<0.
\]
原题要求证明存在非零 $x_0$ 使
\[
f(x_0)=x_0^{\mathsf T}Ax_0.
\]
\end{example}
\exampleblank

\begin{solution}
按原题字面,
\[
f(x)=x^{\mathsf T}Ax
\]
本就是定义,因此结论对任意非零 $x_0$ 都成立,命题没有实质内容.

结合本节上下文,原题极可能漏写了右端的 ``$=0$''.若 intended 结论为
\[
f(x_0)=0,\qquad x_0\neq0,
\]
则取连续函数
\[
\phi(t)=f((1-t)x_1+tx_2).
\]
有 $\phi(0)>0,\phi(1)<0$,由介值定理存在 $t_0\in(0,1)$ 使
\[
\phi(t_0)=0.
\]
若相应向量恰为零,则 $x_1,x_2$ 共线,但同一直线上的二次型取值不可能一正一负,
故该向量非零.
\end{solution}


%例题7.2.5
\begin{example}
设 $A$ 为 $n$ 阶实对称矩阵,各阶顺序主子式
\[
\Delta_1,\ldots,\Delta_n
\]
均非零.证明二次型 $x^{\mathsf T}Ax$ 可经非退化线性替换化为
\[
\lambda_1y_1^2+\cdots+\lambda_ny_n^2,
\qquad
\lambda_k=\dfrac{\Delta_k}{\Delta_{k-1}},
\quad
\Delta_0=1.
\]
\end{example}
\exampleblank

\begin{solution}
由于各顺序主子式均非零,$A$ 存在不选主元的分解
\[
A=LDL^{\mathsf T},
\]
其中 $L$ 为单位下三角矩阵,
\[
D=\diag(d_1,\ldots,d_n).
\]
取前 $k$ 阶行列式得
\[
\Delta_k=d_1d_2\cdots d_k,
\]
所以
\[
d_k=\dfrac{\Delta_k}{\Delta_{k-1}}.
\]
令
\[
y=L^{\mathsf T}x,
\]
即可得到所需标准形.
\end{solution}


\subsection{零点集合与子空间}

\begin{lemma}[\markstar]
设 $A$ 为实对称矩阵,则
\[
A\text{ 半正定或半负定}
\]
的充分必要条件为:对任意 $x$,
\[
x^{\mathsf T}Ax=0
\Longrightarrow
Ax=0.
\]
\end{lemma}

\begin{proposition}[\markstar]
设
\[
V_0=\{x\in\mathbb{R}^n:x^{\mathsf T}Ax=0\}.
\]
则 $V_0$ 是线性子空间的充分必要条件是二次型半正定或半负定.
此时
\[
\dim V_0=n-r(A).
\]
\end{proposition}


%例题7.2.6
\begin{example}
设 $A^{\mathsf T}=A$,并存在 $x_1,x_2$ 使
\[
x_1^{\mathsf T}Ax_1>0,
\qquad
x_2^{\mathsf T}Ax_2<0.
\]
令
\[
W=\{x\in\mathbb{R}^n:x^{\mathsf T}Ax=0\}.
\]
原题进一步提出:
\begin{enumerate}
    \item 证明 $W\cap\mathbb{R}^n=\{0\}$;
    \item 确定 $W$ 中一个极大线性无关组;
    \item 判断 $W$ 是否为 $\mathbb{R}^n$ 的线性子空间.
\end{enumerate}
\end{example}
\exampleblank

\begin{solution}
第一问按原题字面明显不成立,因为
\[
W\subseteq\mathbb{R}^n,
\]
所以
\[
W\cap\mathbb{R}^n=W.
\]
并且由二次型既取正值又取负值,可知存在非零零点,故 $W\neq\{0\}$.
此处原题存在印刷问题.

对第二问,设规范形的正,负惯性指数分别为 $p,q$,零指数为 $r$,
其中 $p,q>0$.在规范坐标中可取
\[
e_i+f_1\quad(i=1,\ldots,p),
\]
\[
e_1-f_1,\qquad e_1+f_j\quad(j=2,\ldots,q),
\]
再加上全部零方向基向量.上述向量总数恰为 $n$,并且线性无关,
因此 $W$ 中存在含 $n$ 个向量的极大线性无关组.

第三问由前述命题立即得到:因为二次型既取正值又取负值,不可能半正定或半负定,
所以
\[
\boxed{W\text{ 不是线性子空间}.}
\]
\end{solution}


%例题7.2.7
\begin{example}
设实二次型
\[
f(x)=x^{\mathsf T}Ax
\]
秩为 $n$,正,负惯性指数分别为 $p,q$,且
\[
p\ge q>0.
\]
\begin{enumerate}
    \item 证明存在 $q$ 维子空间 $W$,使 $f$ 在 $W$ 上恒为零;
    \item 令
    \[
    T=\{x:f(x)=0\},
    \]
    判断 $T$ 是否等于 $W$.
\end{enumerate}
\end{example}
\exampleblank

\begin{solution}
在规范坐标下
\[
f=
y_1^2+\cdots+y_p^2
-z_1^2-\cdots-z_q^2.
\]
令
\[
W=\Span\{e_1+f_1,\ldots,e_q+f_q\}.
\]
对
\[
w=\sum c_i(e_i+f_i)
\]
有
\[
f(w)=\sum c_i^2-\sum c_i^2=0,
\]
故 $\dim W=q$ 且 $f|_W\equiv0$.

但 $T$ 是整个零锥,一般不是线性子空间,而 $W$ 是其中一个最大维的全零子空间.
因此
\[
\boxed{T\neq W}.
\]
\end{solution}


%例题7.2.8
\begin{example}
设 $A\in M_n(\mathbb{R})$ 为对称矩阵,令
\[
T_1=\{x:x^{\mathsf T}Ax=0\},
\qquad
T_2=\{x:Ax=0\}.
\]
证明
\[
T_1=T_2
\iff
T_1\text{ 是 }\mathbb{R}^n\text{ 的线性子空间}.
\]
\end{example}
\exampleblank

\begin{solution}
若 $T_1=T_2$,则 $T_1=\ker A$,当然是子空间.

反之若 $T_1$ 是子空间,由前述命题可知二次型半正定或半负定.
在半正定或半负定情形,
\[
x^{\mathsf T}Ax=0
\iff
Ax=0.
\]
故
\[
T_1=T_2.
\]
\end{solution}


\subsection{正负惯性指数}

\begin{proposition}[\markstar]
设
\[
l_i=c_{i1}x_1+\cdots+c_{in}x_n,
\]
则二次型
\[
f=l_1^2+\cdots+l_p^2-l_{p+1}^2-\cdots-l_{p+q}^2
\]
的正惯性指数不超过 $p$,负惯性指数不超过 $q$.
\end{proposition}


%例题7.2.9
\begin{example}
设实二次型矩阵 $A,B$ 满足
\[
A=D^{\mathsf T}BD,
\qquad
B=C^{\mathsf T}AC,
\]
其中 $C,D$ 不要求可逆.证明 $A,B$ 具有相同的规范形.
\end{example}
\exampleblank

\begin{solution}
由合同型惯性指数单调性,
\[
A=D^{\mathsf T}BD
\Longrightarrow
p(A)\le p(B),\qquad q(A)\le q(B).
\]
另一方面,
\[
B=C^{\mathsf T}AC
\Longrightarrow
p(B)\le p(A),\qquad q(B)\le q(A).
\]
故正,负惯性指数分别相等.又
\[
r(A)=p(A)+q(A)=r(B),
\]
零指数也相同.因此二者规范形相同.
\end{solution}


%例题7.2.10
\begin{example}[\markstar]
设 $A,B$ 为 $n$ 阶实对称矩阵,并且
\[
A=C^{\mathsf T}BC.
\]
证明
\[
p(A)\le p(B),
\qquad
q(A)\le q(B).
\]
\end{example}
\exampleblank

\begin{solution}
把 $B$ 化成规范形
\[
\diag(I_p,-I_q,O).
\]
于是
\[
x^{\mathsf T}Ax=(Cx)^{\mathsf T}B(Cx)
\]
是至多 $p$ 个线性函数平方之和减去至多 $q$ 个线性函数平方之和.
由前述命题,
\[
p(A)\le p,\qquad q(A)\le q.
\]
即得结论.
\end{solution}


%例题7.2.11
\begin{example}
设 $A,B$ 均为 $n$ 阶半正定矩阵.证明
\[
p(A+B)\le p(A)+p(B).
\]
\end{example}
\exampleblank

\begin{solution}
半正定矩阵的正惯性指数等于秩.并且
\[
\ker(A+B)=\ker A\cap\ker B.
\]
因此
\[
r(A+B)
\le r(A)+r(B),
\]
即
\[
\boxed{p(A+B)\le p(A)+p(B)}.
\]
\end{solution}


%例题7.2.12
\begin{example}
设 $A,B$ 为 $n$ 阶半正定矩阵,且
\[
r(A)\le p,\qquad r(B)\le q.
\]
证明 $A-B$ 的正惯性指数不超过 $p$,负惯性指数不超过 $q$.
\end{example}
\exampleblank

\begin{solution}
由惯性指数的加法估计,
\[
p(A-B)\le p(A)+p(-B)=r(A)\le p,
\]
因为 $-B$ 没有正惯性指数.类似地,
\[
q(A-B)\le q(A)+q(-B)=r(B)\le q.
\]
\end{solution}


\subsection{实对称矩阵的正交对角化}

%例题7.2.13
\begin{example}
将下列实对称矩阵正交对角化,并给出一个正交过渡矩阵:
\begin{enumerate}
    \item
    \[
    A_1=
    \begin{pmatrix}
    1&1&2\\
    1&0&3\\
    2&3&1
    \end{pmatrix};
    \]
    \item
    \[
    A_2=
    \begin{pmatrix}
    1&1&1&1\\
    1&2&2&2\\
    1&2&3&3\\
    1&2&3&4
    \end{pmatrix};
    \]
    \item
    \[
    A_3=
    \begin{pmatrix}
    0&1&1&-1\\
    1&0&-1&1\\
    1&-1&0&1\\
    -1&1&1&0
    \end{pmatrix}.
    \]
\end{enumerate}
\end{example}
\exampleblank

\begin{solution}
\begin{enumerate}
    \item 特征多项式为
    \[
    p(\lambda)
    =
    \lambda^3-2\lambda^2-13\lambda-2.
    \]
    记其三个实根为
    \[
    \lambda_1<\lambda_2<\lambda_3.
    \]
    对任一根 $\lambda$,可取特征向量
    \[
    v(\lambda)=
    \begin{pmatrix}
    2\lambda+3\\
    3\lambda-1\\
    \lambda^2-\lambda-1
    \end{pmatrix}.
    \]
    由于 $A_1$ 为实对称矩阵,属于不同特征值的这些向量彼此正交.
    因此令
    \[
    q_i=\dfrac{v(\lambda_i)}{\|v(\lambda_i)\|},
    \qquad
    Q=(q_1,q_2,q_3),
    \]
    即有
    \[
    \boxed{
    Q^{\mathsf T}A_1Q
    =
    \diag(\lambda_1,\lambda_2,\lambda_3).
    }
    \]

    \item 该矩阵满足
    \[
    (A_2)_{ij}=\min(i,j).
    \]
    令
    \[
    \theta_k=\dfrac{(2k-1)\pi}{9},
    \qquad k=1,2,3,4.
    \]
    则特征值为
    \[
    \boxed{
    \lambda_k=
    \dfrac{1}{2-2\cos\theta_k}.
    }
    \]
    相应单位特征向量可取
    \[
    q_k=
    \dfrac23
    \begin{pmatrix}
    \sin\theta_k\\
    \sin2\theta_k\\
    \sin3\theta_k\\
    \sin4\theta_k
    \end{pmatrix}.
    \]
    令 $Q=(q_1,q_2,q_3,q_4)$,则
    \[
    \boxed{
    Q^{\mathsf T}A_2Q
    =
    \diag(\lambda_1,\lambda_2,\lambda_3,\lambda_4).
    }
    \]

    \item 特征值为
    \[
    1\quad(3\text{ 重}),\qquad -3.
    \]
    可取正交单位特征向量
    \[
    q_1=\dfrac1{\sqrt2}(1,1,0,0)^{\mathsf T},
    \]
    \[
    q_2=\dfrac1{\sqrt6}(1,-1,2,0)^{\mathsf T},
    \]
    \[
    q_3=\dfrac1{2\sqrt3}(1,-1,-1,-3)^{\mathsf T},
    \]
    \[
    q_4=\dfrac12(1,-1,-1,1)^{\mathsf T}.
    \]
    令
    \[
    Q=(q_1,q_2,q_3,q_4),
    \]
    则
    \[
    \boxed{
    Q^{\mathsf T}A_3Q=\diag(1,1,1,-3).
    }
    \]
\end{enumerate}
\end{solution}


%例题7.2.14
\begin{example}
设三阶实对称矩阵 $A$ 的秩为 $2$,并且
\[
AB=C,
\]
其中
\[
B=
\begin{pmatrix}
1&1\\
0&0\\
-1&1
\end{pmatrix},
\qquad
C=
\begin{pmatrix}
-1&1\\
0&0\\
1&1
\end{pmatrix}.
\]
求 $A$ 的全部特征值,相应特征向量,并求 $A$.
\end{example}
\exampleblank

\begin{solution}
设 $B$ 的两列为 $b_1,b_2$,$C$ 的两列为 $c_1,c_2$.由
\[
Ab_1=c_1,\qquad Ab_2=c_2
\]
得到
\[
A(e_1-e_3)=(-1,0,1)^{\mathsf T},
\]
\[
A(e_1+e_3)=(1,0,1)^{\mathsf T}.
\]
相加,相减可得
\[
Ae_1=e_3,\qquad Ae_3=e_1.
\]
由于 $A$ 对称,
\[
A=
\begin{pmatrix}
0&0&1\\
0&a&0\\
1&0&0
\end{pmatrix}.
\]
又 $r(A)=2$,故 $a=0$.因此
\[
\boxed{
A=
\begin{pmatrix}
0&0&1\\
0&0&0\\
1&0&0
\end{pmatrix}.
}
\]
其特征值与一组特征向量为
\[
\lambda=1:\ (1,0,1)^{\mathsf T},
\]
\[
\lambda=-1:\ (1,0,-1)^{\mathsf T},
\]
\[
\lambda=0:\ (0,1,0)^{\mathsf T}.
\]
\end{solution}


%例题7.2.15
\begin{example}
设 $A\in M_n(\mathbb{R})$ 为实对称矩阵,其中心化子为
\[
C(A)=\{X\in M_n(\mathbb{R}):AX=XA\}.
\]
证明
\[
\dim C(A)\ge n,
\]
且等号成立当且仅当 $A$ 有 $n$ 个互不相同的特征值.
\end{example}
\exampleblank

\begin{solution}
正交对角化
\[
Q^{\mathsf T}AQ
=
\diag(
\lambda_1I_{m_1},
\ldots,
\lambda_sI_{m_s}
),
\]
其中 $\lambda_i$ 两两不同,
\[
m_1+\cdots+m_s=n.
\]
与该对角矩阵可交换的矩阵恰好按这些特征子空间分块为
\[
\diag(X_1,\ldots,X_s),
\qquad
X_i\in M_{m_i}(\mathbb{R}).
\]
因此
\[
\dim C(A)=\sum\limits_{i=1}^{s}m_i^2.
\]
由
\[
\displaystyle\sum m_i^2\ge\sum m_i=n,
\]
得
\[
\dim C(A)\ge n.
\]
等号成立当且仅当所有 $m_i=1$,即 $A$ 有 $n$ 个互异特征值.
\end{solution}


\begin{proposition}[\markstar]
设 $A$ 是 $n$ 阶实对称矩阵,特征值按
\[
\lambda_1\le\lambda_2\le\cdots\le\lambda_n
\]
排列,则
\[
\lambda_1
=
\min_{x\neq0}
\dfrac{x^{\mathsf T}Ax}{x^{\mathsf T}x},
\qquad
\lambda_n
=
\max_{x\neq0}
\dfrac{x^{\mathsf T}Ax}{x^{\mathsf T}x}.
\]
\end{proposition}


%例题7.2.16
\begin{example}
设实二次型
\[
f(x)=x^{\mathsf T}Ax,
\]
其中 $A$ 是三阶实对称矩阵,并满足
\[
A^3-6A^2+11A-6I=O.
\]
计算
\[
\max_A\ \min_{\|x\|=1} f(x).
\]
\end{example}
\exampleblank

\begin{solution}
多项式分解为
\[
t^3-6t^2+11t-6
=
(t-1)(t-2)(t-3).
\]
由于 $A$ 为实对称矩阵,其特征值都属于
\[
\{1,2,3\}.
\]
于是
\[
\min_{\|x\|=1}f(x)
=
\lambda_{\min}(A)
\le3.
\]
取
\[
A=3I
\]
时原矩阵方程成立,并且最小特征值为 $3$.故
\[
\boxed{3}.
\]
\end{solution}


%例题7.2.17
\begin{example}
设 $A$ 为 $n$ 阶实方阵,$\lambda=a+bi$ 是 $A$ 的任一复特征值.
\begin{enumerate}
    \item 证明
    \[
    S=\dfrac12(A+A^{\mathsf T})
    \]
    的特征值都为实数;
    \item 若
    \[
    \mu_1\le\cdots\le\mu_n
    \]
    是 $S$ 的特征值,证明
    \[
    \mu_1\le a\le\mu_n;
    \]
    \item 对 $b$ 给出类似估计.
\end{enumerate}
\end{example}
\exampleblank

\begin{solution}
第一问因为 $S$ 为实对称矩阵,所以全部特征值均为实数.

设 $z\in\mathbb{C}^n$ 为 $\lambda$ 的特征向量:
\[
Az=\lambda z.
\]
则
\[
a
=
\Ree\lambda
=
\dfrac{z^*Sz}{z^*z}.
\]
对 Hermitian 矩阵 $S$ 使用 Rayleigh 商估计即得
\[
\mu_1\le a\le\mu_n.
\]

再令
\[
H=\dfrac{A-A^{\mathsf T}}{2i}.
\]
$H$ 是 Hermitian 矩阵.若其特征值为
\[
\nu_1\le\cdots\le\nu_n,
\]
则
\[
b
=
\dfrac{z^*Hz}{z^*z},
\]
所以
\[
\boxed{\nu_1\le b\le\nu_n}.
\]
\end{solution}


%例题7.2.18
\begin{example}
设
\[
M=
\begin{pmatrix}
A&B\\
B^{\mathsf T}&C
\end{pmatrix}
\]
为实对称矩阵,其中 $A$ 为 $s$ 阶,$C$ 为 $t$ 阶.
设 $M$ 的最小,最大特征值分别为 $\lambda_1,\lambda_2$,
而 $A,C$ 的最大特征值分别为 $\mu_1,\mu_2$.证明
\[
\lambda_1+\lambda_2\le\mu_1+\mu_2.
\]
\end{example}
\exampleblank

\begin{solution}
取 $M$ 的单位最大特征向量
\[
\begin{pmatrix}x\\y\end{pmatrix},
\qquad
\|x\|^2+\|y\|^2=1.
\]
若 $x,y$ 都非零,令
\[
u=\dfrac{x}{\|x\|},
\qquad
v=\dfrac{y}{\|y\|}.
\]
考虑 $M$ 在二维子空间
\[
\Span\left\{
\begin{pmatrix}u\\0\end{pmatrix},
\begin{pmatrix}0\\v\end{pmatrix}
\right\}
\]
上的压缩.该二维压缩包含 $\lambda_2$ 这一特征值,另一特征值记为 $\theta$.
由 Rayleigh 商有
\[
\theta\ge\lambda_1.
\]
另一方面压缩矩阵的迹为
\[
u^{\mathsf T}Au+v^{\mathsf T}Cv
\le\mu_1+\mu_2.
\]
故
\[
\lambda_1+\lambda_2
\le
\theta+\lambda_2
=
u^{\mathsf T}Au+v^{\mathsf T}Cv
\le
\mu_1+\mu_2.
\]
若 $x=0$ 或 $y=0$,作同样的一维退化讨论即可.
\end{solution}


\newpage

\section{正定与半正定}

\subsection{平方根}

\begin{proposition}[\marktwostar]
若 $S$ 为实对称正定矩阵,则存在唯一实对称正定矩阵 $S_1$ 使
\[
S=S_1^2.
\]
记
\[
S_1=S^{1/2}.
\]
\end{proposition}


%例题7.3.1
\begin{example}[\markstar]
设 $A\in M_n(\mathbb{R})$ 可逆.证明存在正定矩阵 $P$ 与正交矩阵 $U$,
使
\[
A=PU.
\]
\end{example}
\exampleblank

\begin{solution}
令
\[
P=(AA^{\mathsf T})^{1/2}.
\]
因为 $A$ 可逆,所以 $AA^{\mathsf T}$ 正定,故 $P$ 正定且可逆.
定义
\[
U=P^{-1}A.
\]
则
\[
UU^{\mathsf T}
=
P^{-1}AA^{\mathsf T}P^{-1}
=
P^{-1}P^2P^{-1}
=
I.
\]
故 $U$ 正交,并且
\[
\boxed{A=PU}.
\]
\end{solution}


%例题7.3.2
\begin{example}
设 $A$ 为 $n$ 阶半正定矩阵,$B$ 为 $n$ 阶实方阵.
若存在正整数 $r$ 使
\[
A^rB=BA^r,
\]
证明
\[
AB=BA.
\]
\end{example}
\exampleblank

\begin{solution}
正交对角化
\[
A=Q\diag(\lambda_1,\ldots,\lambda_n)Q^{\mathsf T},
\qquad
\lambda_i\ge0.
\]
于是
\[
A^r
=
Q\diag(\lambda_1^r,\ldots,\lambda_n^r)Q^{\mathsf T}.
\]
在非负实数上,映射 $t\mapsto t^r$ 是单射,因此
\[
\lambda_i^r=\lambda_j^r
\iff
\lambda_i=\lambda_j.
\]
所以 $A$ 与 $A^r$ 的特征空间分解完全相同.与 $A^r$ 可交换的矩阵必保持这些特征空间,
故也与 $A$ 可交换.
\end{solution}


%例题7.3.3
\begin{example}
设 $A,B$ 为实对称半正定矩阵,且
\[
AB+BA=O.
\]
证明
\[
AB=BA=O.
\]
\end{example}
\exampleblank

\begin{solution}
取迹得
\[
2\tr(AB)=0.
\]
而
\[
\tr(AB)
=
\tr(A^{1/2}BA^{1/2}).
\]
矩阵
\[
A^{1/2}BA^{1/2}
\]
半正定,迹为零,因此它只能是零矩阵.于是
\[
B^{1/2}A^{1/2}=O.
\]
从而
\[
BA=O,
\qquad
AB=(BA)^{\mathsf T}=O.
\]
\end{solution}


%例题7.3.4
\begin{example}
设 $A$ 为 $n$ 阶实方阵,$B$ 为 $n$ 阶正定矩阵,且
\[
AB^2=B^2A.
\]
证明
\[
AB=BA.
\]
\end{example}
\exampleblank

\begin{solution}
正交对角化
\[
B=Q\diag(\lambda_1,\ldots,\lambda_n)Q^{\mathsf T},
\qquad
\lambda_i>0.
\]
则
\[
B^2=Q\diag(\lambda_1^2,\ldots,\lambda_n^2)Q^{\mathsf T}.
\]
因 $\lambda_i>0$,
\[
\lambda_i^2=\lambda_j^2
\iff
\lambda_i=\lambda_j.
\]
因此 $B$ 与 $B^2$ 的特征空间分解相同,故二者的中心化子相同.
于是
\[
AB^2=B^2A
\Longrightarrow
AB=BA.
\]
\end{solution}


\subsection{同时对角化}

\begin{proposition}[\markstar]
设 $A,B$ 为 $n$ 阶实对称矩阵,且 $A$ 正定,则存在实可逆矩阵 $P$ 使
\[
P^{\mathsf T}AP=I,
\qquad
P^{\mathsf T}BP
=
\diag(\mu_1,\ldots,\mu_n),
\]
其中 $\mu_i$ 是
\[
\det(\lambda A-B)=0
\]
的全部实根,也就是 $A^{-1}B$ 的特征值.
\end{proposition}


%例题7.3.5
\begin{example}
设 $A,B$ 为 $n$ 阶实对称矩阵,且 $B$ 正定.设
\[
\lambda_1,\ldots,\lambda_n
\]
为
\[
\det(\lambda B-A)=0
\]
的全部根.证明:
\begin{enumerate}
    \item $\lambda_i$ 全为实数;
    \item 存在 $\mathbb{R}^n$ 的一组基
    \[
    \varepsilon_1,\ldots,\varepsilon_n
    \]
    使
    \[
    A\varepsilon_i=\lambda_iB\varepsilon_i,
    \qquad
    \varepsilon_i^{\mathsf T}B\varepsilon_j=\delta_{ij}.
    \]
\end{enumerate}
\end{example}
\exampleblank

\begin{solution}
令
\[
C=B^{-1/2}AB^{-1/2}.
\]
则 $C$ 为实对称矩阵,所以存在正交矩阵 $Q=(q_1,\ldots,q_n)$ 使
\[
Cq_i=\lambda_iq_i,
\]
且全部 $\lambda_i$ 为实数.

令
\[
\varepsilon_i=B^{-1/2}q_i.
\]
则
\[
A\varepsilon_i
=
AB^{-1/2}q_i
=
\lambda_iB^{1/2}q_i
=
\lambda_iB\varepsilon_i.
\]
同时
\[
\varepsilon_i^{\mathsf T}B\varepsilon_j
=
q_i^{\mathsf T}q_j
=
\delta_{ij}.
\]
\end{solution}


%例题7.3.6
\begin{example}
设 $f,g$ 为 $\mathbb{R}^n$ 上的实二次型,并且 $f$ 正定.证明方程组
\[
\begin{cases}
f(x)=1,\\
g(x)=1
\end{cases}
\]
无解的充分必要条件是
\[
f-g
\]
为正定或负定二次型.
\end{example}
\exampleblank

\begin{solution}
若 $f-g$ 正定或负定,则显然不可能同时有 $f(x)=g(x)=1$.

反之设方程组无解.集合
\[
S=\{x:f(x)=1\}
\]
在 $f$ 所定义的内积下是单位球面,且
\[
h(x)=f(x)-g(x)
\]
在 $S$ 上处处非零.由于 $h(-x)=h(x)$,即使 $n=1$ 两点取值符号也一致;
当 $n\ge2$ 时 $S$ 连通,所以 $h$ 在 $S$ 上恒正或恒负.

任意 $x\neq0$ 都可缩放为
\[
y=\dfrac{x}{\sqrt{f(x)}}\in S.
\]
由二次齐次性,
\[
h(x)=f(x)h(y).
\]
故 $h$ 对全部非零 $x$ 恒正或恒负,即 $f-g$ 正定或负定.
\end{solution}


%例题7.3.7
\begin{example}
设 $A,B,A-B$ 都是实对称正定矩阵.证明
\[
B^{-1}-A^{-1}
\]
也是正定矩阵.
\end{example}
\exampleblank

\begin{solution}
由
\[
A-B>0
\]
作合同变换得
\[
B^{-1/2}AB^{-1/2}>I.
\]
正定矩阵取逆会反转 Loewner 次序,所以
\[
(B^{-1/2}AB^{-1/2})^{-1}<I.
\]
而
\[
(B^{-1/2}AB^{-1/2})^{-1}
=
B^{1/2}A^{-1}B^{1/2}.
\]
因此
\[
B^{1/2}(B^{-1}-A^{-1})B^{1/2}>0.
\]
由合同保持正定性,得到
\[
\boxed{B^{-1}-A^{-1}>0}.
\]
\end{solution}


%例题7.3.8
\begin{example}[\markstar]
设 $A,B$ 都是 $n$ 阶正定矩阵.证明
\[
\det(A+B)\ge\det A+\det B.
\]
\end{example}
\exampleblank

\begin{solution}
令
\[
C=A^{-1/2}BA^{-1/2}>0
\]
的特征值为 $\lambda_1,\ldots,\lambda_n>0$.则
\[
\det(A+B)
=
\det A\prod\limits_{i=1}^{n}(1+\lambda_i),
\]
\[
\det B
=
\det A\prod\limits_{i=1}^{n}\lambda_i.
\]
而
\[
\displaystyle\prod\limits_{i=1}^{n}(1+\lambda_i)
\ge
1+\prod\limits_{i=1}^{n}\lambda_i,
\]
因为展开式中其余各项均非负.因此
\[
\boxed{\det(A+B)\ge\det A+\det B}.
\]
\end{solution}


%例题7.3.9
\begin{example}
设 $A\in M_n(\mathbb{R})$ 为实对称矩阵,$B\in M_m(\mathbb{R})$ 正定.
证明存在非零矩阵
\[
H\in M_{m,n}(\mathbb{R})
\]
使
\[
B-HAH^{\mathsf T}
\]
正定.
\end{example}
\exampleblank

\begin{solution}
任取非零
\[
H_0\in M_{m,n}(\mathbb{R}),
\]
并令
\[
H=tH_0.
\]
则
\[
B-HAH^{\mathsf T}
=
B-t^2H_0AH_0^{\mathsf T}.
\]
正定矩阵集合是开集,因此当 $t\neq0$ 且充分小时,上式仍为正定矩阵.
于是所求非零 $H$ 存在.
\end{solution}


%例题7.3.10
\begin{example}
设 $A,B$ 为 $n$ 阶实对称矩阵,$B$ 正定,且
\[
A-B
\]
半正定.
\begin{enumerate}
    \item 若 $\lambda$ 是
    \[
    \det(A-\lambda B)=0
    \]
    的根,证明 $\lambda\ge1$;
    \item 证明
    \[
    \det A\ge\det B.
    \]
\end{enumerate}
\end{example}
\exampleblank

\begin{solution}
令
\[
C=B^{-1/2}AB^{-1/2}.
\]
由 $A-B\ge0$,
\[
C-I\ge0,
\]
故 $C$ 的全部特征值都不小于 $1$.而广义特征值方程
\[
\det(A-\lambda B)=0
\]
的根正是 $C$ 的特征值,所以第一问成立.

进一步
\[
\dfrac{\det A}{\det B}
=
\det C
=
\prod\limits_{i=1}^{n}\lambda_i
\ge1.
\]
因此
\[
\boxed{\det A\ge\det B}.
\]
\end{solution}


\subsection{两个重要不等式}

\begin{proposition}[\markstar]
设 $A,B$ 都是 $n$ 阶半正定矩阵,则
\[
\det(A+B)\ge\det A+\det B
\ge
\max\{\det A,\det B\}.
\]
\end{proposition}


%例题7.3.11
\begin{example}
设 $A,B$ 分别为 $m\times n$ 与 $s\times n$ 的行满秩实矩阵,令
\[
Q=
AB^{\mathsf T}(BB^{\mathsf T})^{-1}BA^{\mathsf T}.
\]
证明:
\begin{enumerate}
    \item
    \[
    AA^{\mathsf T}-Q
    \]
    半正定;
    \item
    \[
    0\le\det Q\le\det(AA^{\mathsf T}).
    \]
\end{enumerate}
\end{example}
\exampleblank

\begin{solution}
令
\[
P=
B^{\mathsf T}(BB^{\mathsf T})^{-1}B.
\]
因为 $B$ 行满秩,$P$ 是到 $B$ 的行空间的正交投影,所以
\[
P^2=P,\qquad P^{\mathsf T}=P,\qquad 0\le P\le I.
\]
于是
\[
Q=APA^{\mathsf T},
\]
并且
\[
AA^{\mathsf T}-Q
=
A(I-P)A^{\mathsf T}\ge0.
\]
这证明第一问.

又因 $A$ 行满秩,
\[
AA^{\mathsf T}>0.
\]
对不等式
\[
0\le Q\le AA^{\mathsf T}
\]
作合同归一化,得到一个特征值均落在 $[0,1]$ 的半正定矩阵,故
\[
0\le
\dfrac{\det Q}{\det(AA^{\mathsf T})}
\le1.
\]
因此
\[
\boxed{
0\le\det Q\le\det(AA^{\mathsf T}).
}
\]
\end{solution}


%例题7.3.12
\begin{example}
设
\[
A=(B,C)
\]
是 $n$ 阶实方阵,其中
\[
B\in M_{n,m}(\mathbb{R}),
\qquad
C\in M_{n,n-m}(\mathbb{R}).
\]
证明
\[
|\det A|^2
\le
\det(B^{\mathsf T}B)\,
\det(C^{\mathsf T}C).
\]
\end{example}
\exampleblank

\begin{solution}
有
\[
|\det A|^2
=
\det(A^{\mathsf T}A)
=
\det
\begin{pmatrix}
B^{\mathsf T}B&B^{\mathsf T}C\\
C^{\mathsf T}B&C^{\mathsf T}C
\end{pmatrix}.
\]
矩阵 $A^{\mathsf T}A$ 半正定.对半正定分块矩阵应用 Fischer 行列式不等式,
\[
\det(A^{\mathsf T}A)
\le
\det(B^{\mathsf T}B)\det(C^{\mathsf T}C).
\]
于是结论成立.
\end{solution}


%例题7.3.13
\begin{example}
设 $S,T$ 均为 $n$ 阶实对称半正定矩阵.证明
\[
\det(S+T)
\ge
\dfrac12
\left(\det S+\det T\right).
\]
\end{example}
\exampleblank

\begin{solution}
事实上前述命题给出了更强的不等式
\[
\det(S+T)\ge\det S+\det T.
\]
因此当然有
\[
\boxed{
\det(S+T)
\ge
\dfrac12(\det S+\det T).
}
\]
\end{solution}


%例题7.3.14
\begin{example}
设 $A$ 为 $n$ 阶实矩阵,
\[
S=\dfrac12(A^{\mathsf T}+A),
\]
并且对任意非零 $x\in\mathbb{R}^n$,
\[
x^{\mathsf T}Ax>0.
\]
证明
\[
\det A\ge\det S,
\]
且等号成立当且仅当 $A$ 是对称矩阵.
\end{example}
\exampleblank

\begin{solution}
因为
\[
x^{\mathsf T}Ax
=
x^{\mathsf T}Sx,
\]
所以 $S$ 正定.令
\[
K=\dfrac12(A-A^{\mathsf T}),
\]
则 $K^{\mathsf T}=-K$,且
\[
A=S+K.
\]
作归一化
\[
A
=
S^{1/2}
\left(
I+S^{-1/2}KS^{-1/2}
\right)
S^{1/2}.
\]
记
\[
C=S^{-1/2}KS^{-1/2}.
\]
$C$ 为实反对称矩阵,其非零特征值成对为
\[
\pm i\mu_j.
\]
因此
\[
\det(I+C)
=
\prod_j(1+\mu_j^2)\ge1.
\]
故
\[
\det A
=
\det S\det(I+C)
\ge
\det S.
\]
等号成立当且仅当全部 $\mu_j=0$,即 $C=0$,也就是 $K=0$.
因此等号成立当且仅当
\[
\boxed{A=A^{\mathsf T}}.
\]
\end{solution}


\begin{lemma}[\markstar]
若 $A=(a_{ij})$ 为正定矩阵,则
\[
\det A
\le
a_{nn}\Delta_{n-1},
\]
其中 $\Delta_{n-1}$ 是左上角 $(n-1)$ 阶顺序主子式.
\end{lemma}

\begin{proposition}[\markstar]
Hadamard 不等式:
\begin{enumerate}
    \item 若 $A=(a_{ij})$ 半正定,则
    \[
    \det A\le a_{11}a_{22}\cdots a_{nn};
    \]
    \item 若 $A=(a_{ij})$ 为任意实方阵,则
    \[
    |\det A|^2
    \le
    \prod\limits_{j=1}^{n}
    \left(
    \sum\limits_{i=1}^{n}a_{ij}^2
    \right).
    \]
\end{enumerate}
\end{proposition}


%例题7.3.15
\begin{example}
设 $A=(a_{ij})$ 为 $n$ 阶实方阵,并满足
\[
|a_{ij}|\le1.
\]
证明
\[
|\det A|^2\le n^n.
\]
\end{example}
\exampleblank

\begin{solution}
第 $j$ 列的平方范数满足
\[
\displaystyle\sum\limits_{i=1}^{n}a_{ij}^2
\le n.
\]
由 Hadamard 不等式,
\[
|\det A|^2
\le
\prod\limits_{j=1}^{n}
\left(
\sum\limits_{i=1}^{n}a_{ij}^2
\right)
\le
n^n.
\]
\end{solution}


\subsection{\texorpdfstring{$AB$ 与 $B^{\mathsf T}AB$}{AB 与 BTAB} 型问题}

\begin{proposition}[\markstar]
设 $A$ 为 $n$ 阶正定矩阵,$B$ 为 $n$ 阶实对称矩阵,则:
\begin{enumerate}
    \item $AB$ 的全部特征值均为实数;
    \item $B$ 正定当且仅当 $AB$ 的全部特征值均大于 $0$;
    \item $B$ 半正定当且仅当 $AB$ 的全部特征值均大于等于 $0$.
\end{enumerate}
\end{proposition}

\begin{proposition}[\markstar]
设 $A$ 为 $n$ 阶实对称矩阵,$B$ 为 $n\times m$ 实矩阵,则:
\begin{enumerate}
    \item 若 $A$ 半正定,则 $B^{\mathsf T}AB$ 半正定;
    \item 若 $A$ 正定,则
    \[
    B^{\mathsf T}AB\text{ 正定}
    \iff
    B\text{ 列满秩}.
    \]
\end{enumerate}
\end{proposition}


%例题7.3.16
\begin{example}
设 $A,B$ 为 $n$ 阶正定矩阵,$A-B$ 也正定,且
\[
AB=BA.
\]
证明
\[
A^2-B^2
\]
正定,并举例说明可交换条件不可缺少.
\end{example}
\exampleblank

\begin{solution}
因为 $A,B$ 都是实对称矩阵且可交换,所以可以同时正交对角化:
\[
Q^{\mathsf T}AQ=\diag(a_1,\ldots,a_n),
\]
\[
Q^{\mathsf T}BQ=\diag(b_1,\ldots,b_n).
\]
由 $A-B>0$ 得
\[
a_i>b_i>0.
\]
因此
\[
a_i^2-b_i^2=(a_i-b_i)(a_i+b_i)>0,
\]
故
\[
A^2-B^2>0.
\]

若去掉可交换条件,可取
\[
A=
\begin{pmatrix}
6&-3\\
-3&3
\end{pmatrix},
\qquad
B=
\begin{pmatrix}
1&-1\\
-1&2
\end{pmatrix}.
\]
则 $A,B,A-B$ 都正定,但
\[
A^2-B^2
=
\begin{pmatrix}
43&-24\\
-24&13
\end{pmatrix},
\]
其行列式为
\[
43\cdot13-24^2=-17<0,
\]
故不是正定矩阵.
\end{solution}


%例题7.3.17
\begin{example}
设 $A,B$ 为 $n$ 阶正定矩阵.证明:
\begin{enumerate}
    \item $AB$ 的全部特征值都大于 $0$;
    \item
    \[
    AB=BA
    \iff
    AB\text{ 为正定矩阵}.
    \]
\end{enumerate}
\end{example}
\exampleblank

\begin{solution}
$AB$ 与
\[
A^{1/2}BA^{1/2}
\]
相似,而后者正定,所以 $AB$ 的全部特征值均为正实数.

若 $AB=BA$,则
\[
(AB)^{\mathsf T}=BA=AB,
\]
所以 $AB$ 对称;结合其正特征值,得到 $AB$ 正定.

反之若 $AB$ 正定,则按定义 $AB$ 对称,于是
\[
AB=(AB)^{\mathsf T}=BA.
\]
\end{solution}


%例题7.3.18
\begin{example}
设 $B_1$ 为实对称矩阵,$B_2$ 为同阶正定矩阵.
证明 $B_1B_2$ 的特征值全为实数,并且 $B_1B_2$ 可对角化.
\end{example}
\exampleblank

\begin{solution}
有
\[
B_1B_2
\sim
B_2^{1/2}B_1B_2^{1/2}.
\]
右端是实对称矩阵,因此有实特征值并且可正交对角化.
相似矩阵具有相同谱以及相同的可对角化性,所以
\[
\boxed{B_1B_2\text{ 在实数域上可对角化,且谱为实数}.}
\]
\end{solution}


%例题7.3.19
\begin{example}
设 $B$ 为 $n$ 阶正定矩阵,
\[
C\in M_{n,m}(\mathbb{R}),
\qquad
n>m,\qquad r(C)=m.
\]
令
\[
A=
\begin{pmatrix}
B&C\\
C^{\mathsf T}&O
\end{pmatrix}.
\]
证明 $A$ 有 $n$ 个正特征值,$m$ 个负特征值.
\end{example}
\exampleblank

\begin{solution}
作分块合同变换,
\[
A
\sim_{\mathrm c}
\begin{pmatrix}
B&O\\
O&-C^{\mathsf T}B^{-1}C
\end{pmatrix}.
\]
因为 $B>0$,所以第一块有 $n$ 个正惯性指数.
又 $C$ 列满秩,所以
\[
C^{\mathsf T}B^{-1}C>0,
\]
于是第二块是 $m$ 阶负定矩阵.
由惯性定理,
\[
\boxed{p(A)=n,\qquad q(A)=m}.
\]
实对称矩阵的惯性指数就是正,负特征值的个数.
\end{solution}


\subsection{综合问题}

%例题7.3.20
\begin{example}
设 $A$ 为 $n$ 阶实对称矩阵
\[
A=
\begin{pmatrix}
b+8&3&3&\cdots&3\\
3&b&1&\cdots&1\\
3&1&b&\cdots&1\\
\vdots&\vdots&\vdots&\ddots&\vdots\\
3&1&1&\cdots&b
\end{pmatrix}.
\]
求实数 $b$ 的范围,使 $A$ 正定.
\end{example}
\exampleblank

\begin{solution}
把矩阵写成分块形式
\[
A=
\begin{pmatrix}
b+8&3\boldsymbol{1}^{\mathsf T}\\
3\boldsymbol{1}&D
\end{pmatrix},
\]
其中 $\boldsymbol{1}\in\mathbb{R}^{n-1}$,
\[
D=(b-1)I_{n-1}+J_{n-1}.
\]
$D$ 的特征值为
\[
b-1\quad(n-2\text{ 重}),
\qquad
b+n-2.
\]
因此 $D>0$ 要求且等价于
\[
b>1.
\]
在此前提下,用 Schur 补:
\[
b+8
-
9\boldsymbol{1}^{\mathsf T}D^{-1}\boldsymbol{1}
>0.
\]
因
\[
D\boldsymbol{1}=(b+n-2)\boldsymbol{1},
\]
故
\[
\boldsymbol{1}^{\mathsf T}D^{-1}\boldsymbol{1}
=
\dfrac{n-1}{b+n-2}.
\]
所以要求
\[
b+8-\dfrac{9(n-1)}{b+n-2}>0.
\]
化简为
\[
\dfrac{(b-1)(b+n+7)}{b+n-2}>0.
\]
结合 $b>1$,最终得到
\[
\boxed{b>1}.
\]
\end{solution}


%例题7.3.21
\begin{example}
设
\[
D=
\begin{pmatrix}
A&C\\
C^{\mathsf T}&B
\end{pmatrix}
\]
为正定矩阵,其中 $A,B$ 分别为 $m$ 阶,$n$ 阶实对称矩阵,
$C$ 为 $m\times n$ 实矩阵.
\begin{enumerate}
    \item 证明 $A$ 可逆;
    \item 证明 $D$ 合同于
    \[
    \begin{pmatrix}
    A&O\\
    O&B-C^{\mathsf T}A^{-1}C
    \end{pmatrix};
    \]
    \item 判断 $C^{\mathsf T}A^{-1}C$ 是否正定.
\end{enumerate}
\end{example}
\exampleblank

\begin{solution}
$A$ 是正定矩阵 $D$ 的主子矩阵,因此
\[
A>0,
\]
从而 $A$ 可逆.

作合同变换
\[
\begin{pmatrix}
I&-A^{-1}C\\
O&I
\end{pmatrix}^{\mathsf T}
D
\begin{pmatrix}
I&-A^{-1}C\\
O&I
\end{pmatrix}
=
\begin{pmatrix}
A&O\\
O&B-C^{\mathsf T}A^{-1}C
\end{pmatrix}.
\]

对任意 $x$,
\[
x^{\mathsf T}C^{\mathsf T}A^{-1}Cx
=
(Cx)^{\mathsf T}A^{-1}(Cx)\ge0,
\]
故它总是半正定.
它正定的充分必要条件是
\[
\ker C=\{0\},
\]
即 $C$ 列满秩.原题条件没有保证这一点,例如 $C=0$ 时它就是零矩阵.
\end{solution}


%例题7.3.22
\begin{example}
设 $A$ 为 $n$ 阶正定矩阵,$\alpha\in\mathbb{R}^n$.
证明:若
\[
\det(A-\alpha\alpha^{\mathsf T})=\det A,
\]
则
\[
\alpha=0.
\]
\end{example}
\exampleblank

\begin{solution}
由矩阵行列式引理,
\[
\det(A-\alpha\alpha^{\mathsf T})
=
\det A
\left(
1-\alpha^{\mathsf T}A^{-1}\alpha
\right).
\]
题设与 $\det A>0$ 给出
\[
\alpha^{\mathsf T}A^{-1}\alpha=0.
\]
因为 $A^{-1}$ 正定,只能有
\[
\boxed{\alpha=0}.
\]
\end{solution}


%例题7.3.23
\begin{example}
若 $A,B$ 为半正定矩阵,证明
\[
\tr(AB)\ge0,
\]
且等号成立当且仅当
\[
AB=O.
\]
\end{example}
\exampleblank

\begin{solution}
有
\[
\tr(AB)
=
\tr(A^{1/2}BA^{1/2}).
\]
矩阵
\[
A^{1/2}BA^{1/2}
\]
半正定,所以迹非负.

若迹为零,则这个半正定矩阵所有特征值都为零,因此
\[
A^{1/2}BA^{1/2}=O.
\]
于是
\[
B^{1/2}A^{1/2}=O,
\]
从而
\[
BA=O,
\qquad
AB=(BA)^{\mathsf T}=O.
\]
反向显然.
\end{solution}


%例题7.3.24
\begin{example}
设 $A$ 为 $n$ 阶半正定矩阵,并分块为
\[
A=
\begin{pmatrix}
B&C\\
C^{\mathsf T}&D
\end{pmatrix},
\]
其中 $B$ 为 $m$ 阶方阵.证明
\[
r(B,C)=r(B).
\]
\end{example}
\exampleblank

\begin{solution}
若
\[
x\in\ker B,
\]
则
\[
\begin{pmatrix}x\\0\end{pmatrix}^{\mathsf T}
A
\begin{pmatrix}x\\0\end{pmatrix}
=
x^{\mathsf T}Bx=0.
\]
因为 $A$ 半正定,
\[
z^{\mathsf T}Az=0
\Longrightarrow
Az=0.
\]
故
\[
A\begin{pmatrix}x\\0\end{pmatrix}=0,
\]
从而
\[
C^{\mathsf T}x=0.
\]
所以
\[
\ker B\subseteq\ker C^{\mathsf T}.
\]
这等价于 $C$ 的每一列都属于 $B$ 的列空间,因此把 $C$ 接在 $B$ 后不会增加秩:
\[
\boxed{r(B,C)=r(B)}.
\]
\end{solution}


%例题7.3.25
\begin{example}
设整数 $n\ge3$,实数 $x_1,\ldots,x_n$ 满足
\[
x_1+\cdots+x_n=0,
\qquad
x_1^2+\cdots+x_n^2=1.
\]
证明
\[
x_1x_2+x_2x_3+\cdots+x_{n-1}x_n+x_nx_1
\le
\cos\dfrac{2\pi}{n}.
\]
\end{example}
\exampleblank

\begin{solution}
令 $C$ 为循环实对称矩阵,其相邻两个位置(包括首尾)元素均为
$\dfrac12$,其余为零.则
\[
x^{\mathsf T}Cx
=
x_1x_2+\cdots+x_nx_1.
\]
$C$ 的特征值为
\[
\cos\dfrac{2k\pi}{n},
\qquad
k=0,1,\ldots,n-1.
\]
特征值 $1$ 对应特征向量
\[
(1,\ldots,1)^{\mathsf T}.
\]
题设
\[
x_1+\cdots+x_n=0
\]
说明 $x$ 位于该特征向量的正交补上.在这一子空间中最大特征值为
\[
\cos\dfrac{2\pi}{n}.
\]
又 $\|x\|=1$,故由 Rayleigh 商
\[
\boxed{
x^{\mathsf T}Cx
\le
\cos\dfrac{2\pi}{n}.
}
\]
\end{solution}


%例题7.3.26
\begin{example}[\markstar]
设 $A$ 为 $n$ 阶实方阵.
\begin{enumerate}
    \item 证明对任意非零 $x\in\mathbb{R}^n$,
    \[
    x^{\mathsf T}Ax>0
    \iff
    A+A^{\mathsf T}\text{ 正定};
    \]
    \item 在上述条件成立时证明
    \[
    \det A>0.
    \]
\end{enumerate}
\end{example}
\exampleblank

\begin{solution}
因为
\[
x^{\mathsf T}A^{\mathsf T}x
=
x^{\mathsf T}Ax,
\]
所以
\[
x^{\mathsf T}(A+A^{\mathsf T})x
=
2x^{\mathsf T}Ax.
\]
第一问成立.

令
\[
S=\dfrac12(A+A^{\mathsf T})>0.
\]
若 $\lambda$ 是 $A$ 的实特征值,对相应实特征向量 $x$ 有
\[
\lambda\|x\|^2=x^{\mathsf T}Ax>0,
\]
所以所有实特征值都为正.
非实特征值成共轭对出现,每一对的乘积为
\[
|\lambda|^2>0.
\]
因此所有特征值之积为正,即
\[
\boxed{\det A>0}.
\]
原题若把第一问写成``对任意 $x$''而未排除 $x=0$,则应补充 $x\neq0$.
\end{solution}


%例题7.3.27
\begin{example}
设 $A_1,\ldots,A_m$ 为 $n$ 阶实对称矩阵,并且
\[
\displaystyle\sum\limits_{i=1}^{m}A_i^2=O.
\]
证明
\[
A_i=O
\qquad
(i=1,\ldots,m).
\]
\end{example}
\exampleblank

\begin{solution}
对任意 $x\in\mathbb{R}^n$,
\[
0
=
x^{\mathsf T}
\left(
\sum\limits_{i=1}^{m}A_i^2
\right)x
=
\sum\limits_{i=1}^{m}
x^{\mathsf T}A_i^2x.
\]
由于 $A_i$ 对称,
\[
x^{\mathsf T}A_i^2x
=
\|A_ix\|^2\ge0.
\]
非负数之和恒为零,故对任意 $i,x$,
\[
A_ix=0.
\]
于是
\[
\boxed{A_i=O}.
\]
\end{solution}
